% 全书主要排版参数；修改以下数值即可统一调整对应元素。
\newcommand{\BookBodyFontSize}{8pt}                 % 正文的字体大小
\newcommand{\BookBodyLineHeight}{12pt}              % 正文的基线间距（行距）
\newcommand{\BookFootnoteFontSize}{6pt}             % 脚注正文的字体大小
\newcommand{\BookFootnoteLineHeight}{8pt}          % 单个多行脚注内部的基线间距
\newcommand{\BookFootnoteNumberWidth}{1.15em}        % 脚注编号栏的固定宽度
\newcommand{\BookFootnoteNumberGap}{.35em}           % 脚注编号与脚注正文之间的水平距离
\newcommand{\BookFootnoteSeparation}{8pt}            % 相邻两个脚注之间的垂直距离
\newcommand{\BookFootnoteBlockGap}{7pt plus 1pt minus .5pt} % 正文与页面底部脚注区之间的弹性距离
\newcommand{\BookChapterTitleFontSize}{16pt}         % 章标题的字体大小
\newcommand{\BookChapterTitleLineHeight}{21pt}       % 多行章标题的基线间距
\newcommand{\BookOpeningNoteFontSize}{6pt}           % 题记／译者案正文的字体大小
\newcommand{\BookOpeningNoteLineHeight}{8pt}        % 题记／译者案正文的基线间距
\newcommand{\BookOpeningNoteLabelFontSize}{8pt}      % “题记／译者案”标签的字体大小
\newcommand{\BookOpeningNoteLabelLineHeight}{10pt}   % “题记／译者案”标签的基线间距
\newcommand{\BookRectoTextCenterShift}{3.5mm}        % 右页视觉中心相对纸张中心向外侧偏移的距离

\geometry{
  paperwidth=125mm,
  paperheight=176mm,
  inner=17mm,
  outer=13mm,
  bindingoffset=3mm,
  top=15mm,
  bottom=18mm,
  headheight=4.5mm,
  headsep=4mm,
  footskip=9mm
}

\AtBeginDocument{%
  \fontsize{\BookBodyFontSize}{\BookBodyLineHeight}\selectfont
  \setlength{\parindent}{2em}%
  \setlength{\parskip}{0pt}%
  \clubpenalty=10000
  \widowpenalty=10000
  \displaywidowpenalty=10000
}

\renewcommand{\footnotelayout}{%
  \fontsize{\BookFootnoteFontSize}{\BookFootnoteLineHeight}\selectfont
  \setlength{\parindent}{0pt}%
  \setlength{\parskip}{0pt}%
}
\setlength{\footnotesep}{\BookFootnoteSeparation}
\setlength{\skip\footins}{\BookFootnoteBlockGap}
% A fixed, right-aligned number column keeps one- and two-digit notes aligned.
\makeatletter
\renewcommand{\@makefntext}[1]{%
  % Our custom number column replaces footmisc's \@makefntext, so invoke
  % \footnotelayout explicitly; otherwise LaTeX keeps the class default
  % \footnotesize and the BookFootnote* controls above appear ineffective.
  \footnotelayout
  \noindent
  \hangindent=\dimexpr\BookFootnoteNumberWidth+\BookFootnoteNumberGap\relax
  \hangafter=1
  \makebox[\BookFootnoteNumberWidth][r]{\@makefnmark}%
  \hspace{\BookFootnoteNumberGap}%
  #1%
}
\makeatother
\renewcommand{\footnoterule}{%
  \kern-3pt
  {\color{BookGray}\hrule width .28\textwidth height .35pt}%
  \kern5pt
}

\renewcommand{\contentsname}{目录}
\renewcommand{\cfttoctitlefont}{\hfill\sffamily\bfseries\fontsize{18pt}{22pt}\selectfont}
\renewcommand{\cftaftertoctitle}{\hfill}
\setlength{\cftbeforetoctitleskip}{0pt}
\setlength{\cftaftertoctitleskip}{18pt}
\setlength{\cftbeforepartskip}{12pt}
\setlength{\cftbeforechapskip}{3pt}
\cftpagenumbersoff{part}
% Preserve the centered Part entries from the user's in-progress TOC edit.
\renewcommand{\cftpartfont}{\hfil\sffamily\large\bfseries\color{HomuraViolet}}
\renewcommand{\cftpartfillnum}[1]{%
  \rightskip=0pt
  \parfillskip=0pt plus 1fil
  \par
}
\renewcommand{\cftchapfont}{\fontsize{9pt}{13pt}\selectfont}
\renewcommand{\cftchappagefont}{\fontsize{9pt}{13pt}\selectfont}
\renewcommand{\cftchapleader}{\cftdotfill{2.8}}

\fancypagestyle{bookbody}{%
  \fancyhf{}
  \fancyhead[LE]{\footnotesize\sffamily\color{BookGray}\BookHeaderMotif\quad\nouppercase{\leftmark}}
  \fancyhead[RO]{\footnotesize\sffamily\color{BookGray}\nouppercase{\rightmark}\quad\BookHeaderMotif}
  \fancyfoot[LE,RO]{\sffamily\small\thepage}
  \renewcommand{\headrulewidth}{0pt}
}
\fancypagestyle{chapteropening}{%
  \fancyhf{}
  \renewcommand{\headrulewidth}{0pt}
}
\pagestyle{bookbody}

\newcommand{\BookChapterKicker}{%
  \ifdefempty{\BookSpecialKicker}{%
    CHAPTER~\ifnum\value{chapter}<10 0\fi\arabic{chapter}%
  }{\BookSpecialKicker}%
}
\newcommand{\BookChapterFormat}{%
  \thispagestyle{chapteropening}%
  \raggedright
  {\sffamily\bfseries\fontsize{8.5pt}{10pt}\selectfont\color{HomuraViolet}%
    \BookChapterKicker\hfill
    {\color{MadokaPink}\runeecho{THE VERY SOIL}}\par}%
  \vspace{5pt}%
  {\color{MadokaPink}\rule{12mm}{1.2pt}\hspace{2mm}%
   \color{HomuraViolet}\rule{20mm}{.4pt}\par}%
  \vspace{8pt}%
  \sffamily\bfseries
  \fontsize{\BookChapterTitleFontSize}{\BookChapterTitleLineHeight}\selectfont
  \color{BookBlack}%
}

\ctexset{
  part = {pagestyle=empty},
  chapter = {
    format=\BookChapterFormat,
    name={},
    number={},
    beforeskip=0pt,
    afterskip=18pt,
    afterindent=true,
    pagestyle=chapteropening
  },
  paragraph = {
    format=\sffamily\bfseries\color{HomuraViolet},
    beforeskip=8pt,
    afterskip=3pt,
    runin=false
  }
}

\renewcommand{\chaptermark}[1]{%
  \markright{第\Chinese{chapter}章\quad #1}%
}

\ifbookprint
  \hypersetup{
    hidelinks,
    unicode=true,
    pdftitle={绝望土上的希冀花},
    pdfauthor={Jed A. Blue；JoinNico}
  }
\else
  \hypersetup{
    unicode=true,
    colorlinks=true,
    linkcolor=HomuraViolet,
    citecolor=HomuraViolet,
    urlcolor=MadokaPink,
    pdftitle={绝望土上的希冀花},
    pdfauthor={Jed A. Blue；JoinNico}
  }
\fi
